% keynote bio
\clearpage
\vspace*{\fill}
%\addcontentsline{toc}{section}{Keynote speaker: Dr.~Victor Lawrence}
\begin{center}
\includegraphics[width=3in]{Screenshot 2026-06-04 094331.png}
\end{center}
\vspace*{\fill}

\textbf{Mr.~Erik Levin} will be retiring after more than three decades of service to education.

I was fortunate to have Mr.~Levin my senior year in S\&E, long ago. We were in amazement of a wise, brilliant engineer who would come in during afternoons and be able to open new frontiers in robotics, electronics, and general nerdery for us. He encouraged us to chase our passions in electrical and mechanical engineering, told us all the secret ways to get motors, and sensors, and sonars for free; how to use an octal D-type latch to avoid blowing up our classic IBM PC that was driving our robot through an actual parallel port. He told us where to find free food in college and where to hang out if we were ever in Massachusetts. He introduced us to Horowitz and Hill. I think I am not alone in my experiences. Students universally whine and brag about being challenged. When they are still here it is more of a whine; but when they leave it is more like a thankful brag about how well Mr.~Levin helped to prepare us for the world beyond S\&E. Mr.~Levin always believes the students can succeed if they put forth the effort, even if they do not yet believe themselves. More than once I remember seeing advanced topics in electrical engineering and computer science and remembering I heard them from Mr.~Levin first. He inspires the students to wonder what is possible, because some things feel impossible until confronted with a little pressure and a little Socratic Levin questioning. It has been a great pleasure to teach alongside him the last two years. 

On behalf of many, many alums, we thank Mr.~Levin for his service, for believing in us, and for the gifts he has given us, opening our minds and destroying the limitations we might have thought were there. Congratulations on a well-earned retirement. We wish you all the best.
\vspace*{\fill}

%Today we celebrate Erik Levin and more than three decades of service to education.
%As I prepared these remarks, I had the opportunity to read messages from former students spanning many graduating classes, as well as a tribute written by Joan Kahn, the founding director of the Science & Engineering Program who hired Erik more than 30 years ago. What struck me was how remarkably consistent their reflections were.
%Students wrote about being challenged. They wrote about long nights solving difficult problems, tracing algorithms, revising projects, and learning that excellence requires persistence. At the time, many of them may not have fully appreciated those expectations. Years later, however, they recognize those experiences as some of the most valuable lessons they received.
%Again and again, former students described arriving at college unusually well prepared. They found themselves succeeding in engineering, computer science, mathematics, research, and countless other fields because Erik had already taught them how to think deeply, work independently, and approach difficult problems with discipline and confidence.
%Joan Kahn wrote that Erik believed students could succeed if they were willing to put forth the effort. He expected them to become independent thinkers and self-learners, and he never underestimated what they were capable of achieving. That belief in students has been a defining characteristic of his career.
%Beyond his expertise in mathematics, engineering, computer science, and science, Erik brought his own unique style and perspective to the classroom. Students always knew exactly who he was, what he valued, and how deeply he cared about their growth. 
%The true measure of a teacher is not found in the courses taught or the years served, but in the lives influenced. The students who wrote these letters—many now engineers, researchers, scientists, and professionals - are part of Erik's legacy. They carry with them lessons that continue to shape their lives long after graduation.
%On behalf of the faculty, staff, students, alumni, and families whose lives you have touched, thank you for your dedication, your passion, your high standards, and your commitment to education.
%Congratulations on a well-earned retirement. We wish you all the best in retirement and many happy years ahead doing the things you enjoy most. 
